%======================================================================
% 中文期刊论文模板（Ayaka Launcher 内置模板库）
% 引擎：XeLaTeX；编译命令：latexmk -xelatex main.tex
% 字体：Fandol（TeX Live 自带，无需安装系统字体）
% 说明：双栏排版（正文五号），标题作者摘要区通栏，符合国内科技期刊
%       通行版式。投稿前请按目标期刊的《征稿简则》调整格式。
%======================================================================
\documentclass[twocolumn,zihao=5,a4paper]{ctexart}

\usepackage{amsmath,amssymb}
\usepackage[a4paper,left=1.8cm,right=1.8cm,top=2.2cm,bottom=2.2cm,columnsep=0.8cm]{geometry}
\usepackage{graphicx}
\usepackage{booktabs}

\begin{document}

%======================================================================
% 标题、作者与摘要（通栏）
%======================================================================
\twocolumn[{%
\begin{minipage}{\textwidth}
  \centering
  {\zihao{2}\heiti 基于注意力机制的图像识别方法研究}\par
  \vspace{0.8em}
  {\zihao{-4} 张\quad 三\textsuperscript{1}，李\quad 四\textsuperscript{2}}\par
  \vspace{0.5em}
  {\zihao{-5}
    （1.\ XX 大学 计算机学院，XX 省 XX 市 000000；\quad
     2.\ XX 理工大学 信息工程学院，XX 省 XX 市 000000）}\par
  \vspace{0.4em}
  \begin{minipage}{0.92\textwidth}
    \zihao{-5}
    {\heiti 摘\quad 要：}针对传统卷积神经网络在复杂场景下识别精度不足的问题，
    提出一种结合通道—空间双重注意力机制的改进网络模型。该模型通过通道注意力筛选
    显著性特征通道，通过空间注意力定位目标关键区域，并以跳跃连接保留原始特征信息。
    在 CIFAR-10 与 ImageNet 子集上的实验结果表明，与基线网络相比，所提方法的
    识别准确率分别提升 1.6 与 1.2 个百分点，且参数量仅增加 3.2\%，验证了该方法在
    精度与开销之间取得了良好平衡。\par
    \vspace{0.3em}
    {\heiti 关键词：}图像识别；卷积神经网络；注意力机制；特征提取\par
    \vspace{0.3em}
    {\heiti 中图分类号：}TP391.4\qquad {\heiti 文献标志码：}A\par
  \end{minipage}
  \par
\end{minipage}
\vspace{0.8em}
}]

%======================================================================
\section{引\quad 言}

\quad 图像识别是计算机视觉的核心任务之一，广泛应用于安防监控、医学影像分析、
工业质检等领域\textsuperscript{[1]}。传统方法依赖人工设计特征与浅层分类器，
在光照变化、遮挡等复杂场景下识别效果有限。近年来，深度学习特别是卷积神经网络
（Convolutional Neural Network，CNN）的发展显著提升了图像识别
精度\textsuperscript{[2-3]}，成为该领域的主流技术路线。

然而，标准 CNN 对图像各区域均匀分配计算资源，在背景干扰复杂、目标尺度变化剧烈
的场景下仍存在识别精度不足的问题。注意力机制通过为不同区域分配自适应权重，
可引导网络聚焦关键信息\textsuperscript{[4-5]}，为解决上述问题提供了新思路。
本文在残差网络基础上引入通道—空间双重注意力模块，主要工作包括：
1）设计轻量化的双重注意力模块并分析其作用机理；
2）在公开数据集上验证方法有效性并开展消融实验。

\section{改进网络模型}

\subsection{总体结构}

\quad 本文以 ResNet-18 为骨干网络，在其第 3、4 个残差阶段之后分别插入双重注意力
模块（Dual Attention Module，DAM），并采用跳跃连接融合注意力加权前后的特征。

\subsection{双重注意力模块}

\quad 设输入特征图为 $X \in \mathbb{R}^{H \times W \times C}$，通道注意力权重为
\begin{equation}
  a_c = \sigma\left( W_2\, \delta\left( W_1\, z \right) \right),
  \quad z = \frac{1}{HW} \sum_{i=1}^{H} \sum_{j=1}^{W} X_{ijc},
\end{equation}
其中 $W_1 \in \mathbb{R}^{C/r \times C}$、$W_2 \in \mathbb{R}^{C \times C/r}$
为全连接层参数，$r$ 为缩减比，$\delta(\cdot)$、$\sigma(\cdot)$ 分别为 ReLU 与
Sigmoid 激活函数。空间注意力权重为
\begin{equation}
  a_s = \sigma\left( f\left( \left[ \mathrm{AvgPool}(X);\, \mathrm{MaxPool}(X) \right] \right) \right),
\end{equation}
其中 $f(\cdot)$ 为 $7 \times 7$ 卷积，$[\cdot\,;\cdot]$ 表示通道拼接。输出特征
$\tilde{X} = a_c \odot a_s \odot X$，$\odot$ 为逐元素乘。

\section{实验与分析}

\subsection{实验设置}

\quad 实验基于 PyTorch 2.2 框架，GPU 为 NVIDIA RTX 4090。训练采用 SGD 优化器，
初始学习率 0.1，批大小 128，训练 200 轮。评价指标为 Top-1 识别准确率。

\subsection{结果与分析}

\quad 不同方法在 CIFAR-10 测试集上的对比结果如表~1 所示。

\begin{table}[htbp]
  \centering
  \zihao{-5}
  \caption{不同方法在 CIFAR-10 上的对比结果}
  \begin{tabular}{lccc}
    \toprule
    方法 & 准确率/\% & 参数量/M & 时延/ms \\
    \midrule
    ResNet-18（基线） & 94.8 & 11.2 & 6.4 \\
    SE-ResNet-18      & 95.5 & 11.8 & 7.0 \\
    本文方法          & 96.4 & 11.6 & 6.9 \\
    \bottomrule
  \end{tabular}
\end{table}

\quad 由表~1 可见，本文方法在准确率与模型开销之间取得了良好平衡。消融实验进一步
表明，通道注意力与空间注意力的性能增益分别为 0.9 与 1.1 个百分点，二者具有互补性。

\section{结\quad 论}

\quad 本文提出了一种结合通道—空间双重注意力机制的图像识别方法，通过轻量化注意力
模块增强网络对关键区域与关键通道的特征提取能力。实验结果表明，所提方法能以较小的
参数量代价取得明显的精度提升。后续工作将研究注意力模块的自动化搜索及其在目标检测
任务中的推广。

%======================================================================
% 参考文献
%======================================================================
\begin{thebibliography}{9}
  \zihao{-5}
  \bibitem{ref1} LeCun Y, Bengio Y, Hinton G. Deep learning[J]. Nature, 2015, 521(7553): 436-444.
  \bibitem{ref2} 何恺明, 张翔宇, 任少卿, 等. 深度残差学习中的图像识别[C]//
    计算机视觉与模式识别会议. 拉斯维加斯: IEEE, 2016: 770-778.
  \bibitem{ref3} 周志华. 机器学习[M]. 北京: 清华大学出版社, 2019.
  \bibitem{ref4} Hu J, Shen L, Sun G. Squeeze-and-excitation networks[C]//
    计算机视觉与模式识别会议. 盐湖城: IEEE, 2018: 7132-7141.
  \bibitem{ref5} 王五, 赵六. 基于注意力机制的细粒度图像分类综述[J]. 计算机应用, 2022, 42(6): 1701-1710.
\end{thebibliography}

\end{document}
